\documentclass[11pt]{article}

\usepackage[margin=1in]{geometry}
\usepackage[T1]{fontenc}
\usepackage[utf8]{inputenc}
\usepackage{lmodern}
\usepackage{microtype}
\usepackage{amsmath,amssymb,amsthm,mathtools}
\usepackage{booktabs}
\usepackage{enumitem}
\usepackage{tabularx}
\usepackage{xcolor}
\usepackage[hidelinks]{hyperref}
\usepackage[nameinlink,capitalise]{cleveref}

\newtheorem{theorem}{Theorem}[section]
\newtheorem{proposition}[theorem]{Proposition}
\newtheorem{corollary}[theorem]{Corollary}
\theoremstyle{definition}
\newtheorem{definition}[theorem]{Definition}
\newtheorem{example}[theorem]{Example}
\theoremstyle{remark}
\newtheorem{remark}[theorem]{Remark}

\newcolumntype{Y}{>{\raggedright\arraybackslash}X}
\newcolumntype{L}[1]{>{\raggedright\arraybackslash}p{#1}}

\newcommand{\cA}{\mathcal A}
\newcommand{\cC}{\mathcal C}
\newcommand{\cD}{\mathcal D}
\newcommand{\cE}{\mathcal E}
\newcommand{\cG}{\mathcal G}
\newcommand{\cM}{\mathcal M}
\newcommand{\cN}{\mathcal N}
\newcommand{\cO}{\mathcal O}
\newcommand{\cP}{\mathcal P}
\newcommand{\cR}{\mathcal R}
\newcommand{\cS}{\mathcal S}
\newcommand{\cT}{\mathcal T}
\newcommand{\RR}{\mathbb R}
\newcommand{\PP}{\mathbb P}
\newcommand{\1}{\mathbf 1}
\newcommand{\dd}{\mathrm d}
\newcommand{\MTT}{Modal Triplet Theory}

\title{\vspace{-1em}
When Is a Configuration Physical?\\[0.25em]
\large Admissibility, Chart Failure, and the Logic of Vacuum Selection}
\author{Peter Nero}
\date{Version 2, July 2026}

\begin{document}
\maketitle

\begin{abstract}
Vacuum selection involves several logically different questions that are
often compressed into one.  A configuration can exist as a mathematical
solution, admit a controlled effective chart, satisfy a chosen set of
stability and consistency tests, be reached with nonzero weight under a
dynamics or measure, and agree with observations.  None of these statements
implies all the others.

This paper gives an admissibility-first account of the problem while
correcting two overly strong inferences.  First, failure of one chart or one
certificate does not prove that the underlying mathematical configuration
does not exist.  Global exclusion requires either a chart-independent
obstruction or a transition-compatible exhaustion of the allowed atlas.
Second, exclusion is not selection.  If two configurations remain
admissible, the admissibility predicate alone supplies neither a unique
outcome nor a probability distribution.  Selection additionally requires a
declared dynamics, initial law, measure, conditioning rule, optimization
principle, or other source of relative weight.

We formalize these distinctions with a typed vacuum-selection record that
contains the candidate space and equivalence relation, admissibility rows,
chart transitions, dynamics or ensemble, initial data, stopping and reset
semantics, observation map, conditioning event, and normalization
certificate.  We also show why disconnected admissible components do not by
themselves imply physical inaccessibility: that conclusion depends on
whether the selected evolution is continuous, permits jumps or tunneling,
or changes charts.

The resulting framework preserves the useful core of the
admissibility-first proposal.  It turns consistency, stability, and
controlled-description tests into a provenance-aware filter before
statistical or dynamical selection is attempted.  Current finite
\MTT{} calculations provide examples of such certificates at explicitly
limited tiers.  They do not yet define a global vacuum measure, prove
exhaustion of all branches, or select one observed universe.
\end{abstract}

\section*{Revision note: Version 2}
\addcontentsline{toc}{section}{Revision note: Version 2}
\begin{description}[leftmargin=2.5cm,style=nextline]
\item[Supersedes]
Version 2 supersedes Version 1.0, DOI
\href{https://doi.org/10.5281/zenodo.18255208}
{10.5281/zenodo.18255208}.

\item[Reason]
Version 1 treated failure of an admissible description as if it established
nonexistence of a physical configuration.  It also described admissibility
as selection by exclusion, inferred physical inaccessibility from
disconnected admissible regions, and attributed the lack of convergence in
landscape research to those claims without supplying an atlas-exhaustion
theorem, dynamics, transition law, measure, observation map, or
normalization.

\item[Resolution]
The revision separates mathematical existence, chart validity, kinematic
admissibility, dynamical or statistical realization, and empirical
identification.  It proves elementary no-go propositions for the two invalid
implications, specifies the data required for a genuine vacuum-selection
statement, and compares this role with landscape statistics, tunneling,
effective field theory, and swampland-style consistency filters.

\item[Retained content]
The central methodological idea survives: one should not compare or weight
candidate configurations before stating which mathematical and physical
structures are controlled on them.  Admissibility can sharply reduce a
candidate set and expose missing assumptions.

\item[Remaining boundary]
The paper does not identify the realized vacuum.  Current \MTT{} results do
not yet provide a global exhaustion theorem, a complete visible--hidden
physical compactification, a universal transition dynamics, a normalized
measure on all allowed branches, or a derivation of observed
four-dimensional data from that selection problem.
\end{description}

\tableofcontents

\section{The question in its corrected form}
\label{sec:question}

The phrase ``Which vacuum is physical?'' can hide several distinct
questions:
\begin{enumerate}[leftmargin=2em]
\item Does a candidate exist as a mathematical solution of the stated
      equations and constraints?
\item Is there a chart or representation in which the quantities used in
      the argument are defined and controlled?
\item Does the candidate satisfy a declared set of consistency, stability,
      and approximation tests?
\item Is it reached, occupied, or weighted by a specified dynamics,
      ensemble, measure, or variational rule?
\item Does an observation map send it to data compatible with experiment?
\end{enumerate}
The old paper treated these questions as different descriptions of one
boundary.  They are instead different layers of a complete inference.

The corrected admissibility-first thesis is modest but useful:
\begin{quote}
Before assigning weights or drawing physical conclusions, state the
candidate space, the chart in which each test is meaningful, and the
certificate that each test passes.  Then state separately what selects
among the surviving candidates.
\end{quote}
This is not a rejection of landscape counting, vacuum transitions,
anthropic conditioning, or effective field theory.  It is an ordering rule.
The relevant objects must be typed before they can be combined.

\subsection{Five meanings of ``physical''}

\begin{table}[htbp]
\centering
\small
\begin{tabularx}{\textwidth}{L{0.18\textwidth}YY}
\toprule
Layer & Question & Typical certificate \\
\midrule
Formal existence
& Does the object solve the declared mathematical equations?
& Exact construction, existence theorem, or certified numerical solution. \\
Chart validity
& Are the coordinates, operators, and approximations used here defined?
& Domain theorem, transition map, operator-domain statement, or error bound. \\
Admissibility
& Does the candidate pass the selected consistency and control rows?
& Signed reserves for topology, anomaly, positivity, spectral, descent,
  stability, and truncation conditions. \\
Realization
& Why is this admissible candidate reached or weighted rather than another?
& Initial law and dynamics, invariant measure, transition rates,
  optimization principle, or explicit selection axiom. \\
Empirical identification
& What observations does the realized candidate predict?
& Observation map, matching conventions, uncertainty budget, and
  held-out comparison. \\
\bottomrule
\end{tabularx}
\caption{Five layers that must not be collapsed into the word
``physical.''}
\label{tab:layers}
\end{table}

A candidate can pass an earlier layer and fail a later one.  A formal
solution may lack a controlled low-energy chart.  An admissible candidate
may have zero weight under a declared cosmological dynamics.  A dynamically
preferred candidate may disagree with observation.  Conversely, failure of
one representation need not imply failure of all representations.

\subsection{What this paper does not assume}

The argument does not assume that every useful Hilbert-space description
has a spectral gap.  Hilbert spaces and gap conditions are different
objects.  Nor does it assume that every effective field theory breaks down
at one universal threshold.  Each model has its own domain, scales, and
error estimates.  Finally, it does not assume that configuration space
comes with a canonical probability measure.  A measure is additional data,
not a consequence of naming a set.

\section{Candidate spaces, charts, and admissibility}
\label{sec:charts}

\subsection{The candidate space}

Let \(\cC\) denote a class of candidate configurations and let
\(\sim\) identify descriptions that are regarded as physically or
mathematically equivalent.  Depending on the problem, an element of
\(\cC\) might be:
\begin{itemize}[leftmargin=2em]
\item a stationary point of a scalar potential;
\item a solution of field equations modulo gauge equivalence;
\item a compactification together with metric, flux, bundle, and source
      data;
\item a quantum state or representation of an observable algebra; or
\item an effective theory with declared cutoff and matching data.
\end{itemize}
These are not interchangeable notions of vacuum.  The type of
\(\cC\) must be stated before selection is discussed.

An atlas is a family of partial descriptions
\[
  \{(U_\alpha,\varphi_\alpha)\}_{\alpha\in I},
  \qquad U_\alpha\subseteq \cC,
\]
with transition maps on overlaps whenever the descriptions are intended to
represent the same underlying object.  Here ``chart'' is used broadly: it
may be a coordinate patch, gauge fixing, perturbative expansion, effective
theory, numerical parametrization, or finite projected model.

\subsection{Admissibility rows}

Within a chart \(U_\alpha\), suppose the declared tests are represented by
signed normalized reserves
\[
  r_{\alpha j}:U_\alpha\longrightarrow \RR,
  \qquad j=1,\ldots,m_\alpha.
\]
The convention is that positive reserve means the corresponding condition
is certified.  The chart-relative admissible set is
\[
  \cA_\alpha
  =
  \bigl\{x\in U_\alpha:
  r_{\alpha j}(x)>0\ \text{for every required }j\bigr\}.
\]
The rows can encode exact algebraic constraints, strict positivity,
operator-domain control, numerical residual bounds, truncation errors, or
other conditions.  Their provenance and logical status matter.  A
numerically positive surrogate is not an exact theorem unless its error
certificate controls the relevant implication.

\begin{definition}[Chart-relative admissibility record]
A chart-relative admissibility record is
\[
  \mathfrak A_\alpha
  =
  \bigl(U_\alpha,\varphi_\alpha,
        \{(r_{\alpha j},s_{\alpha j},\pi_{\alpha j})\}_{j=1}^{m_\alpha}\bigr),
\]
where \(s_{\alpha j}\) states the semantics of row \(j\) and
\(\pi_{\alpha j}\) identifies its proof, computation, assumptions, and
source version.  The record is \emph{complete for a declared claim} only if
every premise needed by that claim appears as a row or an explicit
assumption.
\end{definition}

This definition prevents a common ambiguity.  The failure of a row says
that one declared certificate has failed.  What follows depends on its
semantics.  Failure of an error bound can mean only that the approximation
is no longer certified; it need not mean that the exact solution has
ceased to exist.

\subsection{Local failure and global exclusion}

\begin{proposition}[One-chart failure is not nonexistence]
\label{prop:chart}
Let \(x\in\cC\) and let \(\mathfrak A_\alpha\) be one chart-relative
admissibility record.  Either
\[
  x\notin U_\alpha
  \quad\text{or}\quad
  x\in U_\alpha\setminus\cA_\alpha
\]
is insufficient, by itself, to prove that \(x\) does not exist in
\(\cC\), or that no admissible description of \(x\) exists.
\end{proposition}

\begin{proof}
The first statement is immediate because chart domains are partial:
\(x\notin U_\alpha\) is compatible with \(x\in U_\beta\) for another
chart.  For the second, a failed row can be chart dependent.  It is
consistent with the stated data that another record
\(\mathfrak A_\beta\) contains \(x\) and all of its required rows are
positive.  Therefore global nonexistence or exclusion does not follow
without an additional theorem relating all allowed charts or establishing
a chart-independent obstruction.
\end{proof}

A familiar analogy is a coordinate singularity.  The failure of longitude
at a pole does not remove the pole from the sphere.  In a physical
calculation the situation can be subtler, but the logic is the same:
breakdown of coordinates, perturbation theory, gauge choice, or finite
truncation must not be relabeled as breakdown of the underlying object.

\begin{corollary}[Requirements for global exclusion]
\label{cor:global}
A valid global exclusion statement needs at least one of:
\begin{enumerate}[leftmargin=2em]
\item a chart-independent obstruction evaluated on \(x\);
\item an atlas-exhaustion theorem plus transition-compatible admissibility
      rows showing that every allowed chart excludes \(x\); or
\item a definition of the candidate class under which failure of the row
      is itself logically equivalent to nonexistence.
\end{enumerate}
\end{corollary}

The second route is demanding.  On overlaps \(U_\alpha\cap U_\beta\), one
must know whether the two records test the same invariant statement or
only different sufficient conditions.  If they are merely sufficient,
failure in both charts can still leave the underlying property undecided.

\section{Exclusion is not selection}
\label{sec:selection}

An admissibility filter can reduce a candidate class dramatically.  That
is already valuable.  But a filter does not generally choose one surviving
element.

\begin{proposition}[Admissibility alone does not select]
\label{prop:noselect}
Let
\[
  \cA=\{x\in\cC:P(x)=1\}
\]
be the set selected by an admissibility predicate.  If \(\cA\) contains
two inequivalent elements \(x_1\not\sim x_2\), then \(P\) alone determines
neither a unique realized configuration nor a unique probability measure
on \(\cA/\!\sim\).
\end{proposition}

\begin{proof}
Both \(x_1\) and \(x_2\) satisfy exactly the information supplied by
\(P\), so that information does not distinguish them.  Moreover, the
Dirac measures \(\delta_{[x_1]}\) and \(\delta_{[x_2]}\), as well as every
convex combination
\[
  p\,\delta_{[x_1]}+(1-p)\,\delta_{[x_2]},
  \qquad 0\leq p\leq1,
\]
are normalized probability measures supported on the admissible quotient.
No value of \(p\) is determined by the predicate.
\end{proof}

\begin{quote}
Admissibility answers ``which candidates survive these tests?''  Selection
answers ``why this survivor, or this distribution over survivors?''
\end{quote}

Selection can be supplied in several inequivalent ways:
\begin{itemize}[leftmargin=2em]
\item deterministic evolution from selected initial data;
\item stochastic dynamics with specified transition law;
\item an invariant or equilibrium measure, together with a proof of
      existence and normalization;
\item a variational principle with a unique optimizer;
\item cosmological nucleation and expansion rates;
\item conditioning on an observation or observer model; or
\item an explicit branch axiom.
\end{itemize}
Each option changes the claim.  None is silently contained in the word
``admissible.''

\subsection{Why this distinction matters}

Suppose a finite scan begins with \(10^6\) candidates and exact consistency
conditions leave only three.  The reduction from \(10^6\) to three is a
strong result.  It can expose correlations, rule out broad mechanisms, and
make later calculations feasible.  Yet it has not selected one of the
three.  Reporting ``three admissible candidates remain'' is more rigorous
and more informative than hiding the remaining ambiguity behind a
selection label.

Conversely, a probability distribution can be defined before all
admissibility questions are settled, but then its support and physical
interpretation are conditional.  The order advocated here is therefore a
discipline of inference, not a claim that only one research strategy is
mathematically possible.

\section{A typed vacuum-selection record}
\label{sec:record}

To turn an admissibility filter into a selection model, the missing data
should be visible in one record.

\begin{definition}[Vacuum-selection record]
\label{def:record}
A typed vacuum-selection record is
\[
\mathfrak V=
\bigl(
\cC,\sim,\{\mathfrak A_\alpha\},
\cD,\mu_0,\cT,\cR,
\cO,E_{\rm obs},\cN
\bigr),
\]
where:
\begin{enumerate}[leftmargin=2em]
\item \(\cC\) and \(\sim\) define candidates and equivalence;
\item \(\{\mathfrak A_\alpha\}\) is the admissibility atlas;
\item \(\cD\) is a deterministic generator, stochastic transition law,
      path measure, or declared ensemble;
\item \(\mu_0\) is initial data or an initial law when required;
\item \(\cT\) gives stopping, terminal, or observation-time semantics;
\item \(\cR\) gives post-exit continuation, chart transition, or reset
      semantics;
\item \(\cO\) maps selected states or histories to observables;
\item \(E_{\rm obs}\) is the conditioning or comparison event; and
\item \(\cN\) proves normalization, well-posedness, and the error or
      uncertainty statement needed for the final claim.
\end{enumerate}
\end{definition}

The record is deliberately broad enough to cover deterministic,
stochastic, quantum, and statistical proposals.  It does not pretend that
these proposals are equivalent.  Instead, it forces the author to state
which one is being used.

\subsection{Three common specializations}

\paragraph{Deterministic selection.}
Let \(X_t=\Phi_t(X_0)\) be a well-posed flow.  If \(X_0\) is fixed and
\(\Phi_t(X_0)\) converges to a unique equivalence class in \(\cA\), the
limit can define a selection.  The convergence basin and dependence on
initial data remain part of the theorem.

\paragraph{Stochastic selection.}
Let \(X_t\) have a transition kernel \(K_t(x,\dd y)\).  An outcome law might
take the form
\[
  \nu(B)
  =
  \PP_{\mu_0}
  \bigl(\cO(X_T)\in B\mid E_{\rm obs}\bigr).
\]
This expression is meaningful only if the process, stopping time \(T\),
conditioning event, and denominator
\(\PP_{\mu_0}(E_{\rm obs})>0\) are defined.  A different initial law or
transition kernel can produce a different \(\nu\) on the same admissible
set.

\paragraph{Ensemble or counting selection.}
For a finite or regulated class, one may assign weights \(w(x)\) and define
\[
  \nu([x])
  =
  \frac{w([x])\,\1_{\cA}([x])}
       {\sum_{[y]\in\cC/\sim}w([y])\,\1_{\cA}([y])}.
\]
The denominator, regulator, equivalence classes, and origin of \(w\) are
part of the result.  Admissibility supplies the indicator, not the weight.

\subsection{A release checklist for selection claims}

\begin{table}[htbp]
\centering
\small
\begin{tabularx}{\textwidth}{L{0.24\textwidth}YY}
\toprule
Object & Required question & Failure mode if omitted \\
\midrule
Candidate type
& What exactly counts as a vacuum?
& Stationary points, compactifications, states, and EFTs are conflated. \\
Equivalence
& Which descriptions represent the same candidate?
& Gauge copies or dual presentations are overcounted. \\
Admissibility atlas
& Where is each test valid and how do charts overlap?
& Local certificate failure is promoted to global exclusion. \\
Dynamics or weight
& What distinguishes surviving candidates?
& Exclusion is mislabeled as unique selection. \\
Initial law
& From what state or distribution does evolution begin?
& Outcome weights are underdetermined. \\
Exit semantics
& What happens when a chart or reserve fails?
& A stopping condition is mistaken for a new state. \\
Observation map
& How are model variables compared with data?
& Internal coordinates are mistaken for observables. \\
Normalization and error
& Is the law finite, normalized, stable, and accurate enough?
& Formal weights are treated as probabilities. \\
\bottomrule
\end{tabularx}
\caption{Minimum questions for a genuine vacuum-selection statement.}
\label{tab:release}
\end{table}

\section{Disconnected regions and physical accessibility}
\label{sec:connected}

Version 1 inferred that disconnected admissible regions could not be
related physically.  Topology alone does not support that conclusion.

\begin{proposition}[Disconnectedness is not inaccessibility]
\label{prop:disconnect}
If an admissible subset \(\cA\subset\cC\) is disconnected, then no
continuous path lying entirely in \(\cA\) joins different connected
components.  This does not imply that a selected physical evolution cannot
move between the components.
\end{proposition}

\begin{proof}
The first statement is the definition of disconnected components.  For
the second, consider a continuous-time Markov chain on two states
\(\{a,b\}\) with positive jump rates in both directions.  The state space
is disconnected in its discrete topology, yet the process reaches either
state from the other with positive probability.  More generally, an
evolution may leave \(\cA\), use another chart, tunnel, jump, or apply a
reset map.  Whether any such mechanism is physically allowed is extra
data in \(\cD\) and \(\cR\).
\end{proof}

Even when the microscopic evolution is continuous, an effective
description can represent a transition as a jump after unresolved degrees
of freedom are integrated out.  Conversely, a formal instanton or
nucleation solution does not by itself establish an observationally
relevant transition rate.  Coleman--De Luccia vacuum decay is an example
of a theory in which transition data are computed from more than the
topology of the vacuum set \cite{ColemanDeLuccia1980}.

\subsection{What a boundary can establish}

An admissibility boundary can establish that a particular certificate or
chart has reached its declared limit.  It can support a first-exit problem
after a generator is supplied.  It does not independently decide among:
\begin{enumerate}[leftmargin=2em]
\item termination of the underlying history;
\item continuation in another chart;
\item deterministic reset to a new state;
\item stochastic reset according to a kernel;
\item tunneling through a region not represented by the chart; or
\item failure only of the approximation while the exact system remains
      regular.
\end{enumerate}
These alternatives are developed at the process level in the
capacity-gated and selection-front papers
\cite{CapacityGated2026,SelectionFronts2026}.  The present paper owns the
logical distinction as it applies specifically to vacuum selection.

\section{Relation to landscape and consistency programs}
\label{sec:landscape}

\subsection{Landscape counting is not one uniform assumption}

The string-vacuum literature contains several different projects.
Bousso and Polchinski showed how quantized four-form fluxes can generate a
dense discretuum and included a membrane-nucleation dynamics
\cite{BoussoPolchinski2000}.  Douglas framed systematic vacuum
classification in terms of ensembles of effective Lagrangians
\cite{Douglas2003}.  Denef and Douglas computed distributions and counts
for classes of flux vacua \cite{DenefDouglas2004}.  These works do not
simply assume that every formal point is equally realized.  They specify
particular candidate classes, constraints, approximations, and weights.

The admissibility-first framework should therefore not be advertised as an
explanation of why landscape research ``failed to converge.''  It instead
offers a bookkeeping discipline that can be applied inside such programs:
which constraints define the counted class, where are the approximations
valid, what supplies the weight, and what observational conditioning is
used?

\subsection{Consistency filters and the swampland}

The landscape/swampland distinction emphasizes that apparently consistent
low-energy theories need not admit ultraviolet completion in quantum
gravity \cite{Vafa2005}.  This is close in logical form to an admissibility
filter: both distinguish a broad formal class from a smaller class passing
additional consistency conditions.  The criteria are not automatically
the same.  An \MTT{} reserve row becomes relevant to a string or quantum
gravity claim only after a theorem identifies its domain and proves the
corresponding physical implication.

The comparison yields a useful division:
\begin{itemize}[leftmargin=2em]
\item a \emph{consistency filter} asks whether a candidate belongs to the
      intended theory;
\item a \emph{control filter} asks whether a chosen approximation or chart
      is reliable;
\item a \emph{stability filter} asks whether specified perturbations remain
      controlled; and
\item a \emph{selection law} assigns realization or observational weight
      among the survivors.
\end{itemize}
A single condition can play more than one role, but the implication must be
proved rather than inferred from terminology.

\subsection{Anthropic conditioning}

Anthropic arguments add an observation or observer condition rather than
an admissibility theorem.  Weinberg's cosmological-constant bound, for
example, conditions on the formation of gravitationally bound structures
\cite{Weinberg1987}.  One may accept or reject a particular conditioning
model, but its logical role is clear: it supplies \(E_{\rm obs}\) in
\cref{def:record}.  It does not follow from positivity of a local reserve.

\section{Effective descriptions and chart boundaries}
\label{sec:eft}

Effective field theory provides a concrete reason to distinguish an object
from one description of it.  An EFT is specified with degrees of freedom,
symmetries, a cutoff or expansion regime, matching data, and an error
organization.  Its breakdown can indicate that omitted operators,
thresholds, or new degrees of freedom matter.  It need not imply that the
underlying theory or state is nonexistent.  The modern EFT viewpoint grew
from precisely this attention to scale and controlled expansions
\cite{WeinbergEFT1979,WilsonKogut1974}.

In the notation of \cref{sec:charts}, an EFT can be represented as one
chart \(U_\alpha\) with rows for scale separation, truncation error,
positivity, and other model-specific conditions.  A neighboring EFT or a
more complete description may provide \(U_\beta\).  To continue a physical
prediction across their overlap one needs:
\begin{enumerate}[leftmargin=2em]
\item a matching map between variables and observables;
\item overlapping error control;
\item agreement to the declared order or norm; and
\item a statement about which quantities are invariant under the change
      of description.
\end{enumerate}
This is the kind of transition-compatible theorem required by
\cref{cor:global}.  Calling every EFT breakdown ``inadmissible physics''
would erase exactly the multiscale structure that makes EFT useful.

\section{The MTT implementation and its present limit}
\label{sec:mtt}

\subsection{What MTT contributes}

\MTT{} organizes a physical claim around a finite ledger of typed
conditions.  In the corrected Foundation, the ledger distinguishes formal
data, domain conditions, control reserves, exit semantics, and claim
status \cite{MTTFoundation2026}.  The normalized capacity paper makes
otherwise incomparable inequalities explicit by recording a slack, scale,
norm, provenance, and active-row information for each condition
\cite{CapacityMargin2026}.

For vacuum questions, this contributes three practical tools:
\begin{enumerate}[leftmargin=2em]
\item \emph{provenance}: every admissibility row points to the theorem or
      computation that supports it;
\item \emph{typing}: topology, anomaly cancellation, spectral control,
      finite projection, and numerical residuals are not treated as one
      scalar fact; and
\item \emph{anti-promotion boundaries}: a finite or rank-restricted
      certificate remains at that tier until a transfer theorem is proved.
\end{enumerate}

The heterotic-vacuum paper applies this distinction directly: finite ansatz
calculations, attraction in a chosen reduced flow, and physical selection
are separate claims \cite{MTTHeteroticSelection2026}.  The
Hull--Strominger paper similarly keeps its fixed-point correspondence
conditional on the geometric and analytic data named in its hypotheses
\cite{MTTStrominger2026}.

\subsection{Current finite certificates}

Current repository results include:
\begin{itemize}[leftmargin=2em]
\item an exact arithmetic selection of \(q=79\) on a declared branch;
\item a literal finite rank-two Cech witness with its cocycle checks;
\item a certified finite projected HYM solution; and
\item a rank-two weighted-theta/Wiener contraction certificate.
\end{itemize}
These are substantive examples of exact, topological, finite-projection,
and analytic admissibility rows.  They show how a candidate branch can
accumulate auditable certificates instead of relying on one qualitative
label.

They do not yet prove:
\begin{itemize}[leftmargin=2em]
\item exhaustion of every allowed global branch;
\item a unique physical rank-three visible--hidden bundle pair;
\item a complete physical Hull--Strominger or worldsheet endpoint;
\item a normalized measure or transition dynamics on all candidate
      compactifications;
\item a unique realized vacuum; or
\item the observed four-dimensional theory from vacuum selection alone.
\end{itemize}
The reproducibility section below identifies the frozen repository commit
and result manifests.  Their role in this paper is contextual: they
illustrate the structure of certified admissibility, not a completed
vacuum-selection theorem.

\subsection{The theorem still needed}

A strong \MTT{} vacuum-selection result would require a chain of the form
\[
\begin{aligned}
\text{candidate atlas}
&\longrightarrow
\text{transition-compatible admissibility rows}\\
&\longrightarrow
\text{exhaustion or declared candidate boundary}\\
&\longrightarrow
\text{selected dynamics or normalized measure}\\
&\longrightarrow
\text{unique class or outcome law}\\
&\longrightarrow
\text{four-dimensional observation map}.
\end{aligned}
\]
Every arrow needs a domain and a verifier.  The current results occupy
important parts of the second line.  They do not yet supply the entire
chain.

\section{A worked logical example}
\label{sec:example}

Consider three inequivalent candidate compactifications
\(\cC/\!\sim=\{v_1,v_2,v_3\}\).  Suppose exact topology and anomaly tests
exclude \(v_3\), while certified analytic rows retain \(v_1\) and \(v_2\):
\[
  \cA/\!\sim=\{v_1,v_2\}.
\]
This establishes a two-thirds reduction and an exact exclusion of
\(v_3\).  It does not choose between \(v_1\) and \(v_2\).

Now add a continuous-time transition generator on the two survivors,
\[
  Q=
  \begin{pmatrix}
    -k_{12} & k_{12}\\
    k_{21} & -k_{21}
  \end{pmatrix},
  \qquad k_{12},k_{21}>0.
\]
Its stationary law is
\[
  \pi_1=\frac{k_{21}}{k_{12}+k_{21}},
  \qquad
  \pi_2=\frac{k_{12}}{k_{12}+k_{21}}.
\]
The admissibility calculation determines the support
\(\{v_1,v_2\}\).  The rates determine the stationary weights.  An
observation map may then condition those weights further.

This elementary example displays the full separation:
\begin{enumerate}[leftmargin=2em]
\item changing the admissibility rows can change the support;
\item changing the rates can change the weights without changing support;
\item changing the initial law matters before stationarity;
\item changing the observation condition can change the posterior law; and
\item none of these changes alters the mathematical existence of the three
      original candidates.
\end{enumerate}
The same distinctions remain necessary in an infinite-dimensional or
quantum model, where proving that each object exists is much harder.

\section{What would count as progress}
\label{sec:progress}

The framework gives several meaningful intermediate achievements.  A
research program need not leap directly to a unique universe.

\begin{description}[leftmargin=3.2cm,style=nextline]
\item[Certified exclusion]
Prove that a named candidate or family violates a chart-independent
condition.

\item[Controlled reduction]
Reduce a finite or regulated candidate class while reporting all surviving
classes and the completeness of the scan.

\item[Atlas transfer]
Prove that the relevant rows and observables agree on chart overlaps, so a
certificate is not tied to one presentation.

\item[Conditional dynamics]
Given a declared generator and initial law, compute transition, capture, or
stationary probabilities with error control.

\item[Selection theorem]
Prove uniqueness, or derive a normalized outcome law, on an exhausted
candidate class.

\item[Empirical closure]
Map the selected result to observables with conventions, renormalization
and threshold transport, uncertainties, and held-out tests.
\end{description}

Each item is independently valuable.  The labels also prevent an exact
finite certificate from being advertised as empirical closure.

\section{Falsifiability and completion contract}
\label{sec:completion}

The admissibility-first proposal becomes scientifically stronger when it
states how it could fail or remain incomplete.

\subsection{Failure modes}

The proposal would fail as a useful vacuum-selection framework if:
\begin{enumerate}[leftmargin=2em]
\item its rows cannot be made invariant or transition compatible across
      the descriptions required by the physical problem;
\item its filters merely restate known equations without producing new
      exclusions, error control, or computational compression;
\item the candidate class cannot be bounded or exhausted enough for the
      selection claim being made;
\item no dynamics, measure, or variational rule can be sourced without
      importing the desired outcome;
\item different equally admissible row choices produce incompatible
      predictions with no higher-level criterion; or
\item the eventual observation map fails held-out empirical tests.
\end{enumerate}

\subsection{Completion certificate}

A future claim of completed vacuum selection should publish:
\begin{enumerate}[leftmargin=2em]
\item the candidate class and equivalence relation;
\item the atlas and transition maps;
\item every admissibility row with proof status and source hash;
\item an exhaustion statement or an explicit limitation of scope;
\item the dynamics, measure, or optimization principle;
\item initial, boundary, stopping, and reset data;
\item the resulting unique class or normalized outcome law;
\item the observation and matching map;
\item an uncertainty and approximation budget; and
\item independent verification that observed values were not used as
      hidden construction inputs where prediction is claimed.
\end{enumerate}
Without items 4--7, the result is a filter rather than a completed
selection theorem.  Without items 8--10, it is not empirical closure.

\section{Conclusion}
\label{sec:conclusion}

The useful insight in an admissibility-first approach is not that
non-admissible configurations cease to exist in every mathematical sense.
It is that physical inference should not outrun its domain certificates.
A formal solution, a controlled chart, an admissible candidate, a realized
history, and an observed world are different objects.

Two corrections follow.  Failure of one chart cannot establish global
nonexistence without an invariant obstruction or atlas-exhaustion theorem.
And admissibility cannot uniquely select when more than one inequivalent
candidate survives.  Dynamics, measures, conditioning, or another source
of relative weight must enter explicitly.

With these corrections, admissibility remains a promising organizing
principle.  It can turn vague appeals to consistency or stability into
auditable rows, prevent finite calculations from being promoted beyond
their domains, and show exactly which object is missing from a proposed
selection chain.  Current \MTT{} results provide nontrivial examples of
this method.  They narrow and certify parts of a candidate branch.  The
global measure, branch exhaustion, physical compactification, and
empirical observation map remain separate research targets.

\section*{Reproducibility statement}
\addcontentsline{toc}{section}{Reproducibility statement}

The no-go propositions in this paper are elementary consequences of the
typed definitions and require no numerical input.  The finite
\MTT{} examples cited as contextual evidence are maintained in the public
results repository identified below, with immutable commit and manifest
hashes.  No measured Standard Model or cosmological value is used to prove
the logical results of this paper.

\bibliographystyle{plain}

% BEGIN MTT MANAGED COMPUTATIONAL EVIDENCE
\section*{Computational Evidence and Reproducibility}
The q79 arithmetic theorem and audit, literal finite rank-two Cech witness, and rank-two Wiener-contraction certificate are contextual examples of finite, topological, and analytic admissibility certificates. They do not construct a complete physical visible-hidden Hull-Strominger vacuum, select a unique realized branch, supply a global measure or dynamics, or derive observed four-dimensional physics.

The referenced rows are frozen to the curated results repository state identified below. Hashes are grouped in eight-character blocks for line breaking.
\begin{quote}\small
\textbf{Repository:} \url{https://github.com/PeterNero/mtt-results-repro}\\
\textbf{Commit:} \texttt{31247ebb 5c22f3fb b5443024 365433c6 ee0bff4a}\\
\textbf{Manifest:} \path{release/result_manifest.json}\\
\textbf{Manifest SHA-256:}\\
\texttt{fb399689 60b00584 631dbf53 1a708e18}\\
\texttt{ef928d6b 6d935119 c185d7f6 32b1e7cd}
\end{quote}

Tier labels are quoted verbatim from that manifest. A row used directly supports only the specific computational statement identified above; a corpus-state cross-check does not prove this paper's local theorems; and an open row is evidence of an unresolved obligation, never of closure.

\par\begingroup\dimen0=\pagegoal\advance\dimen0 by -\pagetotal\ifdim\dimen0<7\baselineskip\newpage\fi\endgroup
\paragraph{Corpus-state cross-checks.}
\begin{itemize}
\setlength{\itemsep}{0.25em}
\item \path{A19/hym_wiener_contraction} (\emph{numeric certified}).\par Weighted-theta Fourier-tail and Wiener contraction certificate.
\item \path{A07/literal_cech_witness} (\emph{derived exact}).\par Literal 81-entry, 729-cocycle finite Cech witness.
\item \path{A11/q79_exact_audit} (\emph{derived exact}).\par Executable q=79 exact-branch audit.
\item \path{A11/q79_exact_theorem} (\emph{derived exact}).\par CRT q=79 theorem on the selected exact branch.
\end{itemize}

No imported row changes theorem ownership or promotes a neighboring claim: all local statements retain their stated hypotheses, domains, and limitations.
% END MTT MANAGED COMPUTATIONAL EVIDENCE

\bibliography{main}

\end{document}
